\section{Instruction Data Landscape：谁写 Prompt，谁写 Completion}

建立训练目标后，下一步是决定 demonstration 从哪里来。课堂没有简单地把数据分成“好/坏”，而是沿 model-generated 到 human-written 的连续谱比较成本、可扩展性、过滤负担与风格继承。

\subsection{一条样本的基本结构}

上一章已经区分三类监督信号，本节先把 SFT 的一条 demonstration 拆开，回答数据来源究竟在哪些位置进入模型。读后面两页时，应分别追踪 prompt 与 completion 的生产者，而不是只评价最终答案是否流畅。SFT 最小单位通常是 instruction/context 与 completion。这个定义看似简单，却隐藏两个关键选择：prompt 是否来自真实用户分布，completion 是否由人写、强模型写或弱模型自举；这两个来源决定了能力覆盖和偏差传递。

\lecturefigure{slide-07.jpg}{Instruction demonstration 的 instance/completion 结构}{官方幻灯片 p.7}

\begin{knowledgebox}{读图：先分 input provenance 与 output provenance}
图中蓝色块是 instruction/context，绿色块是 demonstration。两者可以来自不同生产者：例如人写 prompt、GPT-4 写 answer，或模型同时生成双方。判断数据质量时不能只读 answer 流畅度，还要检查 prompt 是否覆盖目标用户任务。
\end{knowledgebox}

\lecturefigure{slide-08.jpg}{SFT 数据从 model-generated 到 human-written 的连续谱}{官方幻灯片 p.8}

\begin{knowledgebox}{读图：横轴不是单纯的质量排名}
左侧方法便宜、可扩展，却更依赖 seed、teacher model 与过滤；右侧方法昂贵，但能引入真实人类意图和简洁表达。横轴展示 authorship，不保证越靠右越好，也不保证 synthetic data 必然低质。
\end{knowledgebox}

\teachervoice{Rajani 反复强调，Surge、OpenAssistant、Dolly 与 Self-Instruct 的差别首先是“谁写 input、谁写 output”。她不是在做 dataset leaderboard，而是在建立 provenance taxonomy。}

\subsection{三种 Synthetic Data 机制}

同属 model-generated data，Self-Instruct、UltraChat 与 CAMEL 的 human role 并不相同。把三者并排看，可以理解 synthetic pipeline 的核心变量是 seed、retrieval/material、agent roles 与 filtering，而不是一个“用 GPT 生成”开关。

\lecturefigure{slide-09.jpg}{Self-Instruct：seed tasks、bootstrapping 与 filtering}{官方幻灯片 p.9}

\begin{knowledgebox}{读图：扩展链路与误差链路是同一条链}
少量人工 seed 经过 few-shot generation 扩展，再经过 task classification 与过滤进入 task pool。扩展能力来自 teacher model；模式坍缩、重复和错误也从同一 teacher 传播，因此 filtering 不是可选清洁步骤，而是生成算法的一部分。
\end{knowledgebox}

\lecturefigure{slide-10.jpg}{UltraChat：human-in-the-loop 选择材料并迭代生成}{官方幻灯片 p.10}

\begin{knowledgebox}{读图：人不一定写最终答案，也能改变数据分布}
人类检索主题与支撑材料、决定任务和追问方向，强模型生成对话内容。相比纯模型自举，这种机制把人类劳动放在主题选择与迭代精炼上，用较低写作成本换取更明确的覆盖控制。
\end{knowledgebox}

\lecturefigure{slide-11.jpg}{CAMEL：user/assistant role-playing 生成多轮对话}{官方幻灯片 p.11}

\begin{knowledgebox}{读图：两个 agent 的互动仍由初始任务约束}
一个模型扮演 AI assistant，另一个扮演 user，二者围绕高层 task 自行展开。它适合大规模生成 interaction traces，但角色 prompt、termination 和模型共同偏差会决定对话多样性，不能把“多 agent”自动等同于“更真实用户”。
\end{knowledgebox}

\teachervoice{课堂把三种方法排成“人参与程度”不同的生产线：Self-Instruct 重 seed/filter，UltraChat 让人提供材料并反复精炼，CAMEL 只让人定义高层任务。讲者借此说明，synthetic data 不是一种数据，而是一族 workflow。}

\subsection{从 Landscape 到 H4 的选择}

当时公开数据迅速增长，团队仍决定购买全人工 instruction demonstrations。这个决定不是否定 synthetic data，而是为了得到一个可控的人类基线，再用后续实验比较 task mix、length 与 quantity。

\lecturefigure{slide-12.jpg}{2023 年 instruction dataset landscape 的更新版}{官方幻灯片 p.12}

\begin{knowledgebox}{读图：数据集名称需要机制解释}
Self-Instruct/Unnatural Instructions 以模型扩展为主；T0/FLAN 将已有 NLP tasks 转成 instruction format；UltraChat/ShareGPT 依赖模型或用户对话；Dolly/OpenAssistant/Surge 更靠近人工写作。真正可迁移的知识是机制与 provenance，而不是背名字。
\end{knowledgebox}

\lecturefigure{slide-13.jpg}{H4 在 SFT data continuum 中选择 Surge-Instruct}{官方幻灯片 p.13}

\begin{importantbox}{为什么仍然购买人工数据}
团队希望把“human-written”作为实验控制组：若 synthetic teacher 的风格、事实错误或长度偏好已经进入数据，就很难判断下游模型提升来自任务覆盖还是 teacher imitation。人工数据昂贵，却能提供不同的误差结构。
\end{importantbox}

\teachervoice{讲者直白地说 humans inefficient and expensive，但也不能低估 manually created data 的质量。H4 在 2023 年初选择“非常 manual”的路线，是预算约束下的研究设计，不是对所有项目的永久建议。}

\begin{knowledgebox}{术语消化：常见数据源}
\begin{tabularx}{\textwidth}{@{}L{0.18\textwidth}L{0.25\textwidth}X@{}}
\toprule
名称 & 解决的问题 & 本讲中的机制位置 \\
\midrule
Self-Instruct & 从少量 seed 扩展 instruction set & 模型生成 + 分类过滤 \\
UltraChat & 大规模多轮对话 & 人选主题/材料 + 强模型生成 \\
CAMEL & 角色化交互轨迹 & 两个 LLM role-playing \\
FLAN/T0 & 统一多任务 instruction format & 任务重写与 mixture \\
LIMA & 少量高质量示范 & human-written、数量小 \\
OpenAssistant 与 Dolly & 社区或员工人工对话 & human-authored baseline \\
Surge-Instruct & H4 可控人工 demonstrations & vendor workforce + 指定分布 \\
\bottomrule
\end{tabularx}
\end{knowledgebox}

\subsection{本章小结}

Instruction data 的第一性问题是 authorship 与 workflow：谁决定任务、谁产生内容、谁过滤。Synthetic 与 human data 各自携带不同成本和偏差，后续实验必须显式记录这种 provenance。

